\chapter{格点方法的概览}

LQCD 计算是利用费曼路径积分方法对场论的一种非微扰实现。计算的进行方式与分析地求解场论时完全相同，只要我们有能力完成这些计算。出发点是欧几里得时空中的配分函数
\begin{equation}
  Z = \int \mathcal{D}A_\mu\, \mathcal{D}\bar{\psi}\, \mathcal{D}\psi\, e^{-S}
  \label{eq:partition}
\end{equation}
其中 $S$ 是 QCD 作用量
\begin{equation}
  S = \int d^4x \left( \frac{1}{4} F_{\mu\nu}F^{\mu\nu} - \bar{\psi} M \psi \right).
  \label{eq:qcd_action}
\end{equation}
其中 $M$ 是狄拉克算符。费米子由格拉斯曼变量 $\psi$ 和 $\bar{\psi}$ 表示。它们可以被精确地积出，结果为
\begin{equation}
  Z = \int \mathcal{D}A_\mu\, \det M\, e^{\int d^4x \left(-\frac{1}{4} F_{\mu\nu}F^{\mu\nu}\right)}.
  \label{eq:partition_det}
\end{equation}
费米子贡献现在包含在高度非局域的项 $\det M$ 中，配分函数只是对背景规范组态的积分。在对费米子积分之后，可以将作用量写为 $S = S_\mathrm{gauge} + S_\mathrm{quarks} = \int d^4x ( \frac{1}{4} F_{\mu\nu}F^{\mu\nu} - i \sum_i \log(\mathrm{Det}\,M_i) )$，其中求和遍及各种夸克味道，它们由裸夸克质量的值来区分。

由于「淬火」近似（QQCD）在模拟中所起的核心作用，从一开始就定义它是有利的。它由设定 $\mathrm{Det}\,M = \mathrm{constant}$ 组成，这相当于从 QCD 真空中去除真空极化效应。该近似的细节和一些后果将在第 16.8 节讨论。

物理可观测量的结果通过计算期望值获得
\begin{equation}
  \langle O \rangle = \frac{1}{Z} \int \mathcal{D}A_\mu\, O\, e^{-S}.
  \label{eq:expectation}
\end{equation}
其中 $O$ 是以规范场和夸克场的时间有序乘积表示的任意给定算符组合。（路径积分方法中的期望值对应于时间有序关联函数。）实际上，$O$ 中的夸克场通过维克的收缩场定理重新用夸克传播子表示。这样，所有把夸克作为动力学场的依赖都被去除了。费米子量的基本构成单元是费曼传播子，
\begin{equation}
  S_F(y, j, b; x, i, a) = \left( M^{-1} \right)^{y,j,b}_{x,i,a},
  \label{eq:feynman_propagator}
\end{equation}
其中 $M^{-1}$ 是在给定背景场上计算的狄拉克算符的逆。该矩阵的某一元素 $\left( M^{-1} \right)^{y,j,b}_{x,i,a}$ 是夸克从带有自旋-颜色 $i, a$ 的格点 $x$ 传播到格点-自旋-颜色 $y, j, b$ 的振幅。

到目前为止，以上内容都是标准的场论。我们在 QCD 中面临的问题是如何实际计算这些期望值，以及如何从中提取物理可观测量。我将首先以 $\pi$ 介子的质量和衰变常数为例子来说明第二部分。

考虑两点关联函数 $\langle 0| T \left[ \sum_x O_f(x, t) O_i(0, 0) \right] |0\rangle$，其中 $t > 0$，算符 $O$ 选为轴矢流的第四分量 $O_f = O_i = A_4 = \bar{\psi}\gamma_4\gamma_5\psi$，因为它们与 $\pi$ 介子有大的耦合。该两点关联函数给出以下过程的振幅：由「源」算符 $O_i$ 在时空点 $0$ 从真空中产生一个具有 $\pi$ 介子量子数的态；该态经由 QCD 哈密顿量（或更确切地说，将在第 5 节讨论的传递矩阵）演化到点 $(x, t)$；最后在 $(x, t)$ 处被「汇」算符 $O_f$ 湮灭。量子力学规则告诉我们，$O_i$ 将产生一个态，它是所有与 $\pi$ 介子具有相同量子数的哈密顿量本征态的线性组合，即 $\pi$ 介子、$\pi$ 介子的径向激发、处于 $J = 0$ 态的三个 $\pi$ 介子，\ldots。第二条规则是，在欧几里得时间 $t$ 内传播时，能量为 $E$ 的给定本征态获得权重 $e^{-Et}$。因此，两点函数可以写成对所有可能中间态求和的形式
\begin{equation}
  \langle 0| \sum_x O_f(x, t) O_i(0) |0\rangle
  = \sum_n \frac{\langle 0|O_f|n\rangle\langle n|O_i|0\rangle}{2E_n} e^{-E_n t}.
  \label{eq:2pt_spectral}
\end{equation}
为了研究静止 $\pi$ 介子的性质，我们需要将该态从对 $n$ 的求和中分离出来。为此，第一个简化是使用傅里叶投影 $\sum_x$，因为它将态求和限制为零动量态，因此 $E_n \to M_n$。（注意，只需对 $O_i$ 或 $O_f$ 之一进行傅里叶投影即可。）分离 $\pi$ 介子的第二步，即在能量上投影，由两种策略组合而成。其一是巧妙选择算符 $O$，将态求和限制为单一态（理想的选择是将 $O$ 设为 $\pi$ 介子的量子力学波泛函），其二是考察两点函数在大 $t$ 时的行为，此时由于指数衰减，只有与 $O_i$ 耦合的最低能量态的贡献是重要的。于是
\begin{equation}
  \langle 0| \sum_x O_f(x, t) O_i(0) |0\rangle
  \xrightarrow[t\to\infty]{}
  \frac{\langle 0|O_f|\pi\rangle\langle\pi|O_i|0\rangle}{2M_\pi} e^{-M_\pi t}.
  \label{eq:2pt_pion}
\end{equation}
由于 $\langle 0|A_4(p = 0)|\pi\rangle = M_\pi f_\pi$，右边现在是我们想要的两个量的函数。这样，质量和衰变常数就从时间上的指数衰减速率和振幅中提取出来。

现在让我说明式~\eqref{eq:2pt_pion} 的左边如何用我们在路径积分中控制的两种基本量——规范场和夸克传播子——来表示。利用维克收缩，关联函数可以写成两个夸克传播子 $S_F$ 的乘积，
\begin{equation}
  \begin{aligned}
    -&\langle 0| \sum_x \bar{\psi}(x, t)\gamma_4\gamma_5\psi(x, t)\ \bar{\psi}(0, 0)\gamma_4\gamma_5\psi(0, 0)|0\rangle \\
    &\equiv \langle 0| \sum_x S_F(0; x, t)\gamma_4\gamma_5 S_F(x, t; 0)\gamma_4\gamma_5|0\rangle.
  \end{aligned}
  \label{eq:wick_local}
\end{equation}

\begin{figure}[htbp]
  \centering
  \includegraphics[width=0.8\textwidth]{images/fig4.png}
  \caption{$\pi$ 介子两点关联函数的示意图，其中 (A) 为局域插值算符，(B) 为非局域插值算符。}
  \label{fig:4}
\end{figure}

该关联函数如图~\ref{fig:4}A 所示。为了进一步说明这一维克收缩过程，考虑使用规范不变的非局域算符，例如取 $O = \bar{\psi}(x, t)\gamma_4\gamma_5 \left( P e^{i g \int_x^y dz\, A_\mu(z)} \right)\psi(y, t)$，其中 $P$ 表示路径有序。经过维克收缩后，关联函数为
\begin{equation}
  \langle 0| \sum_x S_F(0; x, t)\gamma_4\gamma_5 \left( P e^{i g \int_z^x A_\mu} \right) S_F(z, t; y, 0)\gamma_4\gamma_5 \left( P e^{i g \int_0^y A_\mu} \right)|0\rangle.
  \label{eq:wick_extended}
\end{equation}
它既涉及规范场又涉及夸克传播子。该关联函数将具有与式~\eqref{eq:2pt_spectral} 所示相同的长 $t$ 行为，但振幅将不同，因此它与 $f_\pi$ 的关系将不再简单。改进 $O$ 向 $\pi$ 介子投影的策略背后的想法，是构造这些算符的合适组合以近似 $\pi$ 介子波函数。例如，若 $\phi(x)$ 是 $\pi$ 介子波泛函，则「涂抹」算符 $O = \bar{\psi}(0)\gamma_4\gamma_5 \int d^3x\, \phi(x)\psi(x)$ 是理想的选择。

最后，必须指出的是，为得到对应于物理 $\pi$ 介子传播的两点函数，我们必须按式~\eqref{eq:expectation} 所定义的那样对规范组态求关联函数的平均，即进行泛函积分。

关于如何完成此类期望值计算的概览包含以下步骤。
\begin{itemize}
  \item \textbf{定义有限维系统：} 规范场的杨-米尔斯作用量和费米子的狄拉克算符必须被转写到离散的时空格点上，其方式要保留 QCD 的所有关键性质——规范不变性、手征对称性、拓扑，以及连续场与格点场之间的一一对应关系。这一步最为困难，即使在今天，我们也没有真正令人满意的格点表述能在 $m_q = 0$ 极限下具有手征对称性，并保持连续场与格点场之间的一一对应关系，即没有加倍者。事实上，Nielson-Ninomiya 定理指出，对于平移不变、定域、厄米的格点理论表述，不可能同时具有手征对称性和没有加倍者 [14]。该定理的一个重要后果是，尽管付出了巨大的努力，格点上仍没有实用的手征费米子表述。关于这些问题及解决尝试的综述，参见 [15--17] 和 [18]。
  \item \textbf{离散化误差：} 当用有限差分来近似作用量中的导数时就会遇到这一问题，
    \begin{equation}
      \partial_x \phi(x) \to \Delta_x \phi(x)
      = \frac{1}{2a}\left( \phi(x + a) - \phi(x - a) \right)
      = \frac{1}{2a}\left( e^{a\partial_x} - e^{-a\partial_x} \right)\phi(x).
      \label{eq:derivative}
    \end{equation}
    显然，这会引入正比于格距 $a$ 的平方的离散化误差。这些误差可以通过以下任一方式减小：使用更高阶的差分格式，其系数经过调整以计入重整化效应；或等价地，在作用量中加入适当的无关算符组合，按 $a$ 的阶数逐阶抵消误差。改进规范作用量和狄拉克作用量的各种方法将在第 12 节和第 13 节讨论。在我的讲义中，我将集中讨论狄拉克作用量的三种最常用的离散化——Wilson [19]、Sheikholeslami-Wohlert（三叶草）[20] 和交错 [21]，它们的误差分别为 $O(a)$、$O(\alpha_s a) - O(a^2)$（取决于三叶草项的系数值）和 $O(a^2)$。通过改进作用量和算符来减小离散化误差，提高了结果的可靠性，而对于动力学模拟，由于计算机能力的限制，这甚至可能是必要的。同时，重要的是要注意到，即使使用最简单的表述——Wilson 原始的规范作用量和费米子作用量 [19]，也存在一个明确定义的程序来获得连续极限的结果。不过在这种情况下，有必要通过可靠的外推至 $a = 0$ 来量化和去除离散化误差。
  \item \textbf{生成背景规范场：} 化学势为零时 QCD 的欧几里得作用量 $S = \int d^4x ( \frac{1}{4} F_{\mu\nu}F^{\mu\nu} - \mathrm{Tr}\log M )$ 是实的且有下界。因此，路径积分中的 $e^{-S}$ 类似于统计力学系统配分函数中的玻尔兹曼因子，即它可以被视为用于生成组态的几率权重。由于 $S$ 是一个广延量，$S$ 的取值范围极大，但分布却非常尖锐地集中在使自由能最小的组态附近。靠近峰值的组态主导泛函积分。这样重要的组态可以通过建立马尔可夫链来生成，与例如 Ising 模型的模拟完全类似。它们被称为「重要性抽样」，因为它们是按由权重 $e^{-S}$ 给出的出现几率生成的。关于用于更新组态的方法的讨论，参见 [1] 或 Creutz 和 Sokal 在 [2] 中的讲义。
  \item \textbf{在给定背景规范场上计算夸克传播子：} 对于给定的背景规范组态，费曼夸克传播子 $S_F$ 是一个由三个指标标记的矩阵——格点、自旋和颜色。该矩阵的某一元素给出一个具有某个自旋、颜色和时空点的夸克传播到另一个时空点、自旋和颜色的振幅。在操作上，$S_F$ 就是狄拉克算符 $M$ 的逆。一旦时空被离散化并变成有限大小，狄拉克矩阵也是有限的，其逆可以数值计算。求狄拉克算符的逆是 LQCD 中的关键计算步骤，消耗超过 90\% 的 CPU 周期。构造关联函数只需要从给定源点到格点所有其他位置的 $S_F$。这对应于 $M^{-1}$ 的 12 列（或行），分别对应 12 个自旋和颜色自由度。
  \item \textbf{关联函数：} 关联函数利用算符场的维克收缩表示为夸克传播子和规范场的路径有序乘积。所计算的关联函数有两种——如前所述针对 $\pi$ 介子情况的规范不变关联函数，或是在固定规范下的关联函数。
  \item \textbf{期望值：} 期望值就是对在一组「重要性抽样」组态上计算的关联函数的平均。即使在有限格点上，背景规范组态的集合也是无限的，因为变量是连续取值的。因此，虽然有可能在指定的背景规范组态上计算关联函数，但精确地完成泛函积分并不可行。因此，它使用蒙特卡洛方法数值地进行。
\end{itemize}

对于理解 LQCD 计算的计算方面的最简单方式，我的建议是首先熟悉任意一个简单统计力学系统的数值处理，例如 Ising 模型。差别在于：(i) LQCD 中的自由度要复杂得多——是 SU(3) 链接矩阵而不是 Ising 自旋，并且这些背景组态上的夸克传播子是通过求狄拉克算符的逆来计算的；(ii) 作用量包含高度非局域的项 Ln Det $M$，这使得规范组态的更新非常昂贵；(iii) 关联函数不像比热或磁化率那样是自旋变量的简单乘积，而是链接变量和夸克传播子的复杂函数。

由于 LQCD 是可重整化的量子场论而不是经典统计力学系统，由此产生的微妙之处在以下两方面发挥作用：分析关联函数作为格距 $a$ 的函数的性态时，以及在对输入参数（夸克场和胶子场及其质量）和关联函数中使用的复合算符进行重整化的量子修正中。乍一看，我们似乎已经在 LQCD 中引入了一个额外的参数，即格距 $a$，然而，正如将在第 14 节中展示的，耦合常数 $g$ 和截断 $a$ 并不是独立的量，而是通过重整化群相互联系，即
\begin{equation}
  \Lambda_\mathrm{QCD} = \lim_{a\to 0} \frac{1}{a} e^{-1/(2\beta_0 g^2(a))} \left( \beta_0 g^2(a) \right)^{-\beta_1/(2\beta_0^2)},
  \label{eq:lambdacd}
\end{equation}
其中 $\Lambda_\mathrm{QCD}$ 是 QCD 的非微扰标度，$\beta_0$ 和 $\beta_1$ 是 $\beta$-函数的前两个与方案无关的系数。在统计力学系统中，格距 $a$ 是一个物理量——分子间距离。在场论中，它仅仅是紫外正规化子，最终必须取为零，同时保持物理量（如重整化的耦合常数（$\Lambda_\mathrm{QCD}$ 和夸克质量）、谱等）不变。

格点结果不精确的原因，是因为在数值模拟中我们必须做出若干近似，正如第 16 节所讨论的。相关不确定度的大小由手头的计算能力决定。随着计算机技术、更好的数值算法和更好的理论理解，它们正在稳步改善。在讨论了 LQCD 的表述之后，我将再回到误差问题。
